%titolo{La Paganela}
%autore{}
%album{}
%tonalita{}
%famiglia{Altre}
%gruppo{}
%momenti{}
%identificatore{la_paganela}
%data_revisione{2026_10_28}
%trascrittore{Francesco Endrici}
%video{}
\beginsong{La Paganela}
\beginverse
Voria veder el Trentino
da 'na vista propri bela.
No 'sto a perder masa tempo
e va' su la Paganela.
Paganela, Paganela,
va la su, va la su, va la su,
ma fa priè non perder temp.
\endverse
\beginverse
Cosa èl sta Paganela?
che no sò cosa che l'è?
No che sai cosa che l'è
No che sai cosa che l'è
L'è la zima la più bela,
de più bele no ghe n'è.
L'è la zima la più bela,
de più bele no ghe n'è.
No ghe n'è, no ghe n'è,
no ghe n'è, no ghe n'è, no ghe n'è,
po, po, po.
\endverse
\beginverse
Tote 'nsema 'na putela
e 'na boza, 'na boza de bon vin,
per goder, per goder la Paganela
e la vista, e la vista del Trentin, del Trentin.
Tote 'nsema 'na putela
e 'na bozza de bon vin,
per goder la Paganela
e la vista del Trentin,
per goder la Paganela
e la vista del Trentin.
\endverse
\beginverse
Paganela, Paganela,
o montagna tutta bela.
Paganela, Paganela,
de più bele no ghe n'è.
No ghe n'è, no ghe n'è,
no ghe n'è, no ghe n'è, no ghe n'è.
\endverse
\beginverse
Da là su se vede 'l ziel,
i torenti e le vedrete,
va l'ociada, va 'l pensier
dal confin fino a le strete.
Paganela, Paganela,
va l'ociada, va 'l pensier
da le strette fien 'l confin.
\endverse
\beginverse
Da 'na banda trenta laghi
e d'Asiago l'altipian,
e d'Asiago l'altipian,
e d'Asiago l'altipian.
Da 'na banda trenta laghi
e d'Asiago l'altipiano,
e da l'altra San Martino
e zò zò fin a Milan.
a Milan, a Milan,
a Milan, a Milan, a Milan,
po, po, po.
\endverse
\endsong